\section{Notation: every symbol, defined once}
\label{sec:notation}

Every quantity used in this document is defined here, grouped by theme,
with its units and the section where it first does physical work.  Two
deliberate collisions of astronomical convention are flagged explicitly:
$\beta$ denotes both the ecliptic latitude (sky model,
Sect.~\ref{sec:sky}) and the Moffat wing index (PSF, Sect.~\ref{sec:psf});
$\alpha$ denotes both the lunar phase angle (Sect.~\ref{sec:sky}) and the
Moffat core radius (Sect.~\ref{sec:psf}).  The context — sky or PSF —
always disambiguates them, and no equation mixes the two meanings.

\subsection*{Radiometry and magnitudes (Sect.~\ref{sec:radiometry})}

\begin{tabular}{llp{8.6cm}}
\toprule
Symbol & Units & Definition \\
\midrule
$\lambda$ & \si{\angstrom} & Wavelength; all internal grids are in \si{\angstrom}. \\
$\flam(\lambda)$ & \si{erg.s^{-1}.cm^{-2}.\angstrom^{-1}} & Spectral flux
  density per unit wavelength (``F-lambda''), the quantity all templates
  are converted to. \\
$F_\nu$ & Jy & Spectral flux density per unit frequency;
  $1\,\mathrm{Jy} = 10^{-23}\,\si{erg.s^{-1}.cm^{-2}.Hz^{-1}}$. \\
$Z_\nu$ & Jy & The $F_\nu$ zero point of a magnitude system: the flux
  density of a zero-magnitude source (AB: 3631 Jy at every wavelength;
  Vega: per-filter values from the SVO profile). \\
$m$ & mag & Magnitude: $m = -2.5\log_{10}(F/F_{\rm zero})$ in the stated
  system and band. \\
$m_{v0}$ & mag & Catalogue field: the visual magnitude represented by a
  stored template file as shipped (0 for BPGS spectra). \\
$T_\lambda$ & — & Filter (passband) transmission, $0 \le T_\lambda \le 1$. \\
$W_\lambda$ & — & Synthetic-photometry weighting: $\lambda$ for a
  photon-counting detector, 1 for an energy integrator (SVO
  \code{DetectorType}). \\
$Q_\lambda$ & — & Detector quantum efficiency: detected electrons per
  incident photon at $\lambda$. \\
$O_\lambda$ & — & Optics/instrument throughput curve (everything between
  aperture and detector except the filter and QE). \\
$\eta$ & — & Scalar (grey) efficiency factor for losses with no stated
  wavelength dependence. \\
$A$ & \si{cm^2} & Unobstructed collecting area,
  $A = \tfrac{\pi}{4}(D^2 - d_{\rm obs}^2)$, with $D$ the aperture and
  $d_{\rm obs}$ the central-obstruction diameter. \\
$h,\,c$ & cgs & Planck constant and speed of light; $hc/\lambda$ is the
  photon energy. \\
$\dot S,\ \dot B,\ \dot d$ & \si{e^-/s} & Detected electron rates: source
  (aperture- or element-integrated), sky, and dark current per pixel. \\
$v$ & km/s & Radial velocity of the target; positive = receding. \\
$E(B\!-\!V)$ & mag & Interstellar reddening (colour excess); $R_V =
  A(V)/E(B\!-\!V) = 3.1$ assumed; $A(\lambda)$ = extinction in
  magnitudes at $\lambda$ (CCM89 law). \\
\bottomrule
\end{tabular}

\subsection*{Geometry, time and the atmosphere
(Sects.~\ref{sec:atmosphere}, \ref{sec:sky})}

\begin{tabular}{llp{8.6cm}}
\toprule
Symbol & Units & Definition \\
\midrule
$h$ (altitude) & deg & Apparent altitude above the horizon (includes
  refraction unless stated ``geometric''). \\
$z$ & deg & Zenith distance, $z = 90^\circ - h$. \\
$X$ & — & Airmass: the optical path through the atmosphere relative to the
  zenith path; computed from the apparent altitude with
  Eq.~(\ref{eq:pickering}). \\
$T_{\rm atm}(\lambda)$ & — & Zenith atmospheric transmission; the
  transmission at airmass $X$ is $T_{\rm atm}^{\,X}$ (Bouguer law). \\
$k_\lambda$ & mag/airmass & Extinction coefficient;
  $T_{\rm atm} = 10^{-0.4 k_\lambda}$. \\
$n(\lambda)$ & — & Refractive index of air at the site conditions. \\
$\Delta R(\lambda)$ & arcsec & Differential atmospheric refraction
  (atmospheric dispersion) relative to a reference wavelength,
  Eq.~(\ref{eq:dr}). \\
$H,\ \varphi,\ \delta$ & deg & Hour angle of the target, geographic
  latitude of the site, declination of the target. \\
$q$ & deg & Parallactic angle: the position angle of the local vertical at
  the target, measured from North through East, Eq.~(\ref{eq:parallactic}). \\
$P,\ T_{\rm site},\ E$ & hPa, \si{\celsius}, m & Site pressure,
  temperature (ISA defaults from the elevation) and elevation. \\
$t_{\rm exp}$ & s & Exposure time of a single frame. \\
$d_{\rm cm}$ & cm & Telescope aperture diameter (in the scintillation
  law). \\
$\sigma_{\rm scint}/I$ & — & Fractional rms intensity modulation by
  scintillation, Eq.~(\ref{eq:scint}). \\
$s$, FWHM & arcsec & Seeing: the full width at half maximum of the
  seeing-limited stellar profile. \\
$\sigma_g$ & arcsec & Guiding (tracking) error, rms per axis; adds a blur
  of FWHM $2.355\,\sigma_g$ in quadrature to the seeing. \\
$r_0$ & m & Fried parameter of Kolmogorov turbulence;
  seeing $\propto \lambda/r_0$. \\
\bottomrule
\end{tabular}

\subsection*{Sky brightness (Sect.~\ref{sec:sky})}

\begin{tabular}{llp{8.6cm}}
\toprule
Symbol & Units & Definition \\
\midrule
$\mu$ & mag/arcsec$^2$ & Surface brightness: the magnitude of the flux
  received from one square arcsecond of sky. \\
S10 & — & Surface-brightness unit: the equivalent of one star of
  $m_V = 10$ per square degree;
  $\mu = 27.78 - 2.5\log_{10} q$ mag/arcsec$^2$ for a brightness of $q$
  S10 units. \\
nL & — & nanoLambert, the luminance unit of the Krisciunas--Schaefer
  moonlight model; converted to S10 by the factor 3.8. \\
$q_{\rm air},\,q_{\rm zod},\,q_\ast$ & S10 & Airglow, zodiacal-light and
  integrated-starlight brightness components. \\
$a$ & — & Solar-cycle activity, interpolated between tabulated solar
  minima: 0.8 at minimum, 2.0 at maximum. \\
$V(z)$ & — & van Rhijn airglow path-length factor,
  Eq.~(\ref{eq:vanrhijn}). \\
$\hat h$ & km & Height of the airglow emitting layer (90 km). \\
$R_\oplus$ & m & Mean Earth radius, $6.371\times10^6$ m. \\
$\beta$ (sky) & deg & Ecliptic latitude of the field. \\
$b$ & deg & Galactic latitude of the field. \\
$\epsilon$ & deg & Solar elongation of the field: its angular distance
  from the Sun. \\
$\alpha$ (Moon) & deg & Lunar phase angle: $0^\circ$ = full Moon,
  $90^\circ$ = quarter; computed as $180^\circ$ minus the Sun--Moon
  elongation. \\
$\rho$ & deg & Angular separation between the Moon and the target. \\
$I^\ast$ & — & Lunar illuminance factor of the K\&S model. \\
$f(\rho)$ & — & K\&S scattering function (Rayleigh + aureole),
  Eq.~(\ref{eq:ks}). \\
$X_{\rm m}$ & — & Airmass of the Moon. \\
$d_{\rm moon}$ & km & Topocentric distance of the Moon; the illuminance
  scales as $(\bar d/d_{\rm moon})^2$ with $\bar d = 384\,400$ km. \\
$\mu_{\rm SQM}$ & mag/arcsec$^2$ & Zenith $V$ sky brightness measured by
  a Sky Quality Meter. \\
\bottomrule
\end{tabular}

\subsection*{PSF and coupling (Sect.~\ref{sec:psf})}

\begin{tabular}{llp{8.6cm}}
\toprule
Symbol & Units & Definition \\
\midrule
$\sigma$ & arcsec & Gaussian width, $\sigma = \mathrm{FWHM}/2.355$. \\
$\alpha$ (Moffat) & arcsec & Moffat core radius,
  $\alpha = \mathrm{FWHM}/(2\sqrt{2^{1/\beta}-1})$. \\
$\beta$ (Moffat) & — & Moffat wing index; small $\beta$ = strong wings;
  Gaussian is the $\beta \to \infty$ limit. \\
$\mathrm{EE}(r)$ & — & Encircled energy: fraction of the total PSF flux
  inside a circular aperture of radius $r$. \\
$T(w,\delta)$ & — & Slit coupling: fraction of the PSF flux through an
  infinite slit of width $w$ decentred by $\delta$,
  Eq.~(\ref{eq:slit}). \\
$F_\nu$ (CDF) & — & Student-$t$ cumulative distribution with $\nu$
  degrees of freedom (only inside Eq.~\ref{eq:slit}; not a flux). \\
$p$ & arcsec/pixel & Plate scale: $p = 206265\, s_{\rm pix}/f_{\rm tel}$
  with $s_{\rm pix}$ the pixel size and $f_{\rm tel}$ the telescope focal
  length in the same length units. \\
\bottomrule
\end{tabular}

\subsection*{Noise budget and solvers (Sect.~\ref{sec:photometry})}

\begin{tabular}{llp{8.6cm}}
\toprule
Symbol & Units & Definition \\
\midrule
$S,\ B,\ D$ & e$^-$ & Accumulated source, sky and dark charge in the
  aperture over $t_{\rm exp}$. \\
$N_{\rm pix}$ & — & Number of detector pixels summed in the aperture or
  extraction window. \\
$\sigma_R$ & e$^-$ & Read noise per pixel per frame, rms. \\
$g$ & e$^-$/ADU & Conversion gain; the ADC quantization noise is
  $g/\sqrt{12}$ per pixel. \\
$N_{\rm bit}$ & — & ADC bit depth; digital saturation at
  $g\,(2^{N_{\rm bit}}-1)$ electrons. \\
$n_{\rm sky}$ & — & Number of pixels used to \emph{measure} the
  background (sky annulus or slit sky windows); finite $n_{\rm sky}$
  multiplies background variances by $(1 + N_{\rm pix}/n_{\rm sky})$. \\
$\dot b$ & \si{e^-/s}/pixel & Optional extra background rate (detector
  glow, stray light). \\
$C$ & s$^{1/2}$ & Scintillation constant of a configuration:
  $\sigma_{\rm scint}/I = C\,(2 t_{\rm exp})^{-1/2}$. \\
$Q$ & \si{e^-/s} & Total Poisson-like rate in the solver,
  $Q = \dot S + \dot B + \dot D$ (plus the scintillation rate). \\
$\Omega$ & arcsec$^2$ & Angular area of an extended source. \\
$N,\ T$ & —, s & Number of stacked frames and total integration time
  $T = N\,t_{\rm sub}$. \\
\bottomrule
\end{tabular}

\subsection*{Spectroscopy (Sect.~\ref{sec:spectroscopy})}

\begin{tabular}{llp{8.6cm}}
\toprule
Symbol & Units & Definition \\
\midrule
$R$ & — & Resolving power, $R = \lambda/\Delta\lambda$. \\
$\Delta\lambda$ & \si{\angstrom} & Resolution element: the wavelength
  width the instrument cannot subdivide; one output row of the ETC. \\
$\mathcal{D}$ & \si{\angstrom}/pixel & Linear dispersion at the
  detector. \\
$\ell$ & pixel & Line-spread-function FWHM in dispersion-direction
  pixels (slitless: seeing $\oplus$ intrinsic $\oplus$ dispersion
  smear). \\
$n_{\rm mm}$ & mm$^{-1}$ & Grating groove density. \\
$m$ (grating) & — & Diffraction order (1 for the Star Analyser). \\
$L$ & mm & Grating-to-sensor distance of a converging-beam grating. \\
$f_{\rm coll},\ f_{\rm cam}$ & mm & Collimator and camera focal lengths
  of a slit spectrograph. \\
$w_{\rm slit}$ & arcsec & Slit width on the sky. \\
$\delta x$ & \si{\angstrom} & Width of one dispersion pixel,
  $\delta x = \Delta\lambda/(\text{pixels per element})$. \\
$W$ & \si{\angstrom} & FWHM of a spectral line in the Cayrel relation,
  Eq.~(\ref{eq:cayrel}). \\
EW & \si{\angstrom} & Equivalent width: the width of a rectangle of
  continuum height with the same integrated absorption/emission as the
  line. \\
\bottomrule
\end{tabular}
